% ==========================================================
% UNIVERSIDAD ESTATAL PENÍNSULA DE SANTA ELENA
% FACULTAD DE SISTEMAS Y TELECOMUNICACIONES
% CARRERA DE INGENIERÍA DE SOFTWARE
% ASIGNATURA: FRAMEWORK PROGRAMACIÓN WEB
% ==========================================================
% Informe Técnico: Cuadro Comparativo y Evidencias Swagger UI
% Formato: Una captura por página
% ==========================================================

\documentclass[12pt,a4paper]{article}

% ==========================================================
% PAQUETES
% ==========================================================
\usepackage[utf8]{inputenc}
\usepackage[T1]{fontenc}
\usepackage[spanish,es-tabla]{babel}
\usepackage[top=2.2cm,bottom=2.2cm,left=2.2cm,right=2.2cm]{geometry}
\usepackage{graphicx}
\usepackage[table]{xcolor}
\usepackage{titlesec}
\usepackage{tabularx}
\usepackage{booktabs}
\usepackage{array}
\usepackage{float}
\usepackage{hyperref}
\usepackage{fancyhdr}
\usepackage{lastpage}
\usepackage{microtype}
\usepackage{amsmath}
\usepackage{caption}
\usepackage[most]{tcolorbox}

% ==========================================================
% PALETA FORMAL EN NEGRO
% ==========================================================
\definecolor{formalblack}{HTML}{111111}
\definecolor{pureblack}{HTML}{000000}
\definecolor{headerdark}{HTML}{1E293B}
\definecolor{rowgray}{HTML}{F8FAFC}
\definecolor{lightbg}{HTML}{F8FAFC}

% ==========================================================
% HIPERENLACES
% ==========================================================
\hypersetup{
    colorlinks=true,
    linkcolor=pureblack,
    urlcolor=pureblack,
    citecolor=pureblack,
    pdfauthor={Damian Torres - UPSE},
    pdftitle={Informe Tecnico: Cuadro Comparativo y Pruebas en /docs}
}

% ==========================================================
% ESTILO DE SECCIONES FORMAL EN NEGRO
% ==========================================================
\titleformat{\section}
{\color{formalblack}\Large\bfseries}
{\thesection.}
{0.45em}
{}
[{\vspace{1pt}\color{formalblack}\titlerule[1.0pt]}]

\titlespacing*{\section}{0pt}{10pt}{10pt}

% ==========================================================
% ENCABEZADOS Y PIES DE PÁGINA FORMALES
% ==========================================================
\pagestyle{fancy}
\setlength{\headheight}{14pt}
\fancyhf{}
\lhead{\small\textbf{\textcolor{formalblack}{UPSE}} \textcolor{gray}{$\,\vert\,$} \textcolor{formalblack}{Framework Programación Web}}
\rhead{\small\textcolor{formalblack}{SaaS API Multitenant} \textcolor{gray}{$\,\vert\,$} \textbf{\textcolor{formalblack}{Damián Torres}}}
\cfoot{\small\textcolor{formalblack}{Página \thepage\ de \pageref{LastPage}}}
\renewcommand{\headrulewidth}{0.8pt}
\renewcommand{\footrulewidth}{0.3pt}

% ==========================================================
% CAJAS Y TABLAS
% ==========================================================
\tcbset{
    formalbox/.style={
        enhanced,
        colback=lightbg,
        colframe=formalblack,
        boxrule=0.6pt,
        leftrule=3.5pt,
        arc=1mm,
        left=10pt,
        right=10pt,
        top=8pt,
        bottom=8pt,
        boxsep=0pt
    }
}

\renewcommand{\arraystretch}{1.25}
\setlength{\tabcolsep}{6pt}
\newcolumntype{L}[1]{>{\raggedright\arraybackslash}p{#1}}
\newcolumntype{Y}{>{\raggedright\arraybackslash}X}

\captionsetup{
    font={small,rm},
    labelfont={bf,color=formalblack},
    textfont={color=formalblack},
    skip=6pt
}

\graphicspath{{docs_screenshots/}{./}}

% ==========================================================
% DOCUMENTO
% ==========================================================
\begin{document}

% ==========================================================
% PÁGINA 1: PORTADA INSTITUCIONAL FORMAL
% ==========================================================
\begin{titlepage}
\thispagestyle{empty}
\centering
\vspace*{0.8cm}

\IfFileExists{upse.png}{
    \includegraphics[width=3.8cm]{upse.png}\\[0.6cm]
}{
    {\Large \textbf{UPSE}}\\[0.6cm]
}

{\large \textbf{\textcolor{formalblack}{UNIVERSIDAD ESTATAL PEN\'{I}NSULA DE SANTA ELENA}}}\\[0.2cm]
{\normalsize \textbf{Facultad de Sistemas y Telecomunicaciones}}\\[0.06cm]
{\normalsize \textbf{Carrera de Ingenier\'{i}a de Software}}

\vspace{0.8cm}
\noindent\rule{\textwidth}{1.0pt}
\vspace{0.25cm}

{\normalsize \textbf{FRAMEWORK PROGRAMACI\'{O}N WEB}}

\vspace{0.25cm}
\noindent\rule{\textwidth}{1.0pt}
\vspace{0.8cm}

{\Large \textbf{Informe T\'{e}cnico: Cuadro Comparativo Arquitect\'{o}nico y Validaci\'{o}n Pr\'{a}ctica de Endpoints en FastAPI}\par}
\vspace{0.45cm}
{\large \textit{Desarrollo de Endpoints Multitenant, Validaci\'{o}n Pydantic y Pruebas en Swagger UI (\texttt{/docs})}}

\vfill

\begin{tabularx}{\textwidth}{@{}lX@{}}
\textbf{Estudiante:}            & \textbf{Dami\'{a}n Jes\'{u}s Torres Cadena} \\
                                & \texttt{damian.torrescadena0563@upse.edu.ec} \\[0.25cm]
\textbf{Docente:}               & \textbf{Ing. Carlos Mu\~{n}oz} \\[0.2cm]
\textbf{Asignatura:}            & Framework Programaci\'{o}n Web \\[0.2cm]
\textbf{Per\'{i}odo acad\'{e}mico:} & 2026--2 \\[0.2cm]
\textbf{Fecha de presentaci\'{o}n:} & Noviembre de 2026 \\
\end{tabularx}

\vspace{0.8cm}

\end{titlepage}

\clearpage
\pagenumbering{arabic}
\setcounter{page}{1}

% ==========================================================
% PÁGINA 2: CUADRO COMPARATIVO PROPIO
% ==========================================================
\section{Cuadro Comparativo: MVC Clásico vs. MVT (Django) vs. API-First (FastAPI)}

A continuaci\'{o}n se presenta la comparativa arquitect\'{o}nica requerida, contrastando los tres modelos de desarrollo web en funci\'{o}n de sus dimensiones estructurales, desacoplamiento de capas y aislamiento multitenant:

\begin{table}[H]
\centering
\small
\begin{tabularx}{\textwidth}{@{}L{3.0cm}Y L{3.5cm}Y@{}}
\toprule[1.2pt]
\rowcolor{headerdark}
\textcolor{white}{\textbf{Dimensi\'{o}n}} & 
\textcolor{white}{\textbf{MVC Cl\'{a}sico (Spring/Laravel)}} & 
\textcolor{white}{\textbf{MVT (Django)}} & 
\textcolor{white}{\textbf{API-First (FastAPI)}} \\
\midrule[0.8pt]
\textbf{Controlador y Despacho} & 
El \textbf{Controller} intercepta HTTP, invoca el Modelo y selecciona la Vista para renderizado. & 
La \textbf{View} act\'{u}a de controlador; consulta ORM y despacha contexto a plantillas. & 
Los \textbf{Path Operations} act\'{u}an como controladores desacoplados con serializaci\'{o}n JSON directa. \\
\rowcolor{rowgray}
\textbf{Capa de Vista (UI)} & 
\textbf{View}: Plantillas procesadas en servidor (JSP/Blade). Acoplamiento visual fuerte. & 
\textbf{Template}: HTML con DTL/Jinja. Renderizado del lado del servidor (SSR monol\'{i}tico). & 
\textbf{Headless / Desacoplada}: Frontend aut\'{o}nomo (React, Flutter, IoT) que consume JSON puro. \\
\textbf{Validaci\'{o}n y Contrato} & 
Formularios backend espec\'{i}ficos o anotaciones manuales dependientes del framework. & 
\textbf{Forms / ModelForms}: Ligados r\'{i}gidamente al ORM y a la sesi\'{o}n del navegador. & 
\textbf{Pydantic Schemas}: Validaci\'{o}n \textit{fail-fast} en frontera, tipado est\'{a}tico y contrato OpenAPI autogenerado. \\
\rowcolor{rowgray}
\textbf{Concurrencia / Runtime} & 
Tradicionalmente s\'{i}ncrono; modelo \textit{thread-per-request} (bloqueante en I/O). & 
Predominantemente s\'{i}ncrono sobre WSGI (soporte ASGI incremental no nativo). & 
\textbf{As\'{i}ncrono nativo (ASGI)}: Bucle de eventos no bloqueante optimizado para I/O concurrente. \\
\textbf{Aislamiento Multitenant} & 
Resuelto por middlewares, filtros de servlet o tenancy a nivel de esquema en base de datos. & 
Depende de extensiones (\texttt{django-tenants}) o filtrado manual en QuerySets. & 
\textbf{Segmentaci\'{o}n expl\'{i}cita por ruta (\texttt{/api/\{tenant\_id\}})} y \textit{scoped querying} directo. \\
\bottomrule[1.2pt]
\end{tabularx}
\caption{Cuadro comparativo propio entre MVC Cl\'{a}sico, MVT y API-First.}
\label{tab:comparativa}
\end{table}

\begin{tcolorbox}[formalbox]
\textbf{Conclusi\'{o}n Arquitect\'{o}nica Breve:}
La arquitectura \textbf{API-First} adoptada con FastAPI desacopla totalmente la l\'{o}gica del servicio respecto a la capa de presentaci\'{o}n visual. Pydantic act\'{u}a como barrera perimetral tipada garantizando la integridad de datos antes de que ingresen a la capa l\'{o}gica, mientras que el direccionamiento por tenant provee aislamiento estricto de recursos.
\end{tcolorbox}

\clearpage

% ==========================================================
% PÁGINA 3: CAPTURA 1 - VISTA GENERAL EN SWAGGER UI
% ==========================================================
\section{Evidencia 1: Vista General de Endpoints en Swagger UI (\texttt{/docs})}

Exposici\'{o}n interactiva completa de los contratos OpenAPI en Swagger UI. Se aprecian las operaciones implementadas para la soluci\'{o}n SaaS multitienda (\texttt{GET}, \texttt{POST}, \texttt{PUT} y \texttt{DELETE}) junto con los esquemas tipados de Pydantic:

\vspace{0.4cm}

\begin{figure}[H]
\centering
\setlength{\fboxrule}{0.6pt}\setlength{\fboxsep}{0pt}%
\fbox{\includegraphics[width=0.96\textwidth]{01_swagger_overview.png}}
\caption{Vista general de las operaciones RESTful y esquemas OpenAPI en \texttt{/docs}.}
\label{fig:swagger-overview}
\end{figure}

\begin{tcolorbox}[formalbox]
\textbf{Detalle T\'{e}cnico:} La especificaci\'{o}n OpenAPI 3.1 autogenerada refleja con precisi\'{o}n los par\'{a}metros de ruta requeridos (\texttt{tenant\_id}, \texttt{pedido\_id}) y los esquemas de entrada para registrar y actualizar pedidos.
\end{tcolorbox}

\clearpage

% ==========================================================
% PÁGINA 4: CAPTURA 2 - VALIDACIÓN PYDANTIC (HTTP 422)
% ==========================================================
\section{Evidencia 2: Validación Pydantic con HTTP 422}

Verificaci\'{o}n de la regla de validaci\'{o}n defensiva implementada en \texttt{PedidoEntrada} con \texttt{Field(..., gt=0)}. Al enviar una orden con $\text{cantidad} = 0$ y $\text{total} = -25.0$, el motor intercepta la solicitud y devuelve autom\'{a}ticamente \textbf{HTTP 422 Unprocessable Content}:

\vspace{0.4cm}

\begin{figure}[H]
\centering
\setlength{\fboxrule}{0.6pt}\setlength{\fboxsep}{0pt}%
\fbox{\includegraphics[width=0.96\textwidth]{02_pydantic_validation_422.png}}
\caption{Respuesta HTTP 422 impidiendo registrar pedidos con cantidad $\le 0$ o total $\le 0$.}
\label{fig:pydantic-422}
\end{figure}

\begin{tcolorbox}[formalbox]
\textbf{Detalle T\'{e}cnico:} El cuerpo de respuesta JSON expone detalladamente cada transgresi\'{o}n (\texttt{"type": "greater\_than"}, \texttt{"loc": ["body", "cantidad"]}, \texttt{"msg": "Input should be greater than 0"}), demostrando validaci\'{o}n estricta en frontera.
\end{tcolorbox}

\clearpage

% ==========================================================
% PÁGINA 5: CAPTURA 3 - ENDPOINT PUT (CAMBIO DE ESTADO)
% ==========================================================
\section{Evidencia 3: Endpoint \texttt{PUT} --- Cambio de Estado de Pedido}

Ejecuci\'{o}n del endpoint \texttt{PUT /api/\{tenant\_id\}/pedidos/\{pedido\_id\}/estado}. Se procesa la actualizaci\'{o}n del pedido con $\text{id} = 2$ en \texttt{tienda\_norte}, conmutando su estado a \texttt{"Entregado"} con respuesta exitosa \textbf{HTTP 200 OK}:

\vspace{0.4cm}

\begin{figure}[H]
\centering
\setlength{\fboxrule}{0.6pt}\setlength{\fboxsep}{0pt}%
\fbox{\includegraphics[width=0.96\textwidth]{03_put_estado_200.png}}
\caption{Ejecuci\'{o}n exitosa del endpoint PUT para cambio de estado retornando HTTP 200.}
\label{fig:put-estado}
\end{figure}

\begin{tcolorbox}[formalbox]
\textbf{Detalle T\'{e}cnico:} La mutaci\'{o}n se realiza de manera segura y at\'{o}mica; el servicio valida la existencia previa del tenant y asegura que el pedido pertenezca exclusivamente a la tienda indicada antes de alterar su estado.
\end{tcolorbox}

\clearpage

% ==========================================================
% PÁGINA 6: CAPTURA 4 - ENDPOINT DELETE CON AISLAMIENTO
% ==========================================================
\section{Evidencia 4: Endpoint \texttt{DELETE} --- Aislamiento Multitenant}

Ejecuci\'{o}n del endpoint \texttt{DELETE /api/\{tenant\_id\}/pedidos/\{pedido\_id\}}. Se elimina el pedido con $\text{id} = 1$ perteneciente a \texttt{tienda\_norte} con respuesta \textbf{HTTP 200 OK}, garantizando aislamiento entre tiendas:

\vspace{0.4cm}

\begin{figure}[H]
\centering
\setlength{\fboxrule}{0.6pt}\setlength{\fboxsep}{0pt}%
\fbox{\includegraphics[width=0.96\textwidth]{04_delete_pedido_200.png}}
\caption{Eliminaci\'{o}n de pedido confirmada con HTTP 200 y aislamiento verificado entre tiendas.}
\label{fig:delete-aislamiento}
\end{figure}

\begin{tcolorbox}[formalbox]
\textbf{Garant\'{i}a de Aislamiento Multitenant:} La eliminaci\'{o}n se acota \'{u}nicamente a la partici\'{o}n del tenant especificado en la URL. El pedido con $\text{id} = 1$ perteneciente a \texttt{tienda\_sur} permaneci\'{o} \'{i}ntegro e inalterado en memoria, evitando cualquier interferencia cruzada.
\end{tcolorbox}

\end{document}
